% Generated by build_manuscript.py from 06_conclusions.md
\section{6 Conclusions}\label{conclusions}

This study formulated partial C-scan acquisition as a sequential decision problem for compression-after-impact assessment. A surface prior, a state-dependent acquisition actor and a separate frozen CAI predictor organize the flow from currently visible evidence to the next internal observation. The common-predictor design permits comparisons of acquisition policies without changing the assessment mapping, while the trajectory objective accounts for both prediction quality and the stage at which observations become available.

On 50 selected validation specimens, the main policy achieved 44.286 MPa MAE under a 25\% acquisition cap, compared with 46.910 MPa for geometry-spread, and a lower trajectory area than the shared static and open-loop policies. The signed timing decomposition recovered these area differences from saved events and located the largest main-versus-open-loop contribution difference in the second completion stage. Component results were mixed: no-VLM feedback performed better, all four reported component-related intervals crossed zero, and partial acquisition did not reach the 41.690 MPa complete-input MAE.

The resulting contribution is a reproducible way to connect observation allocation with a mechanical assessment objective and to report its quality--acquisition trade-offs. The observed comparisons apply to offline image replay with validation-based checkpoint selection and one policy-seed panel. Instrument-linked costs and independent specimen evaluation remain necessary before interpreting these image-fraction results as physical inspection-efficiency gains.
